Fix a score vector outside the \(P_X^{\otimes n}\)-null set on which either a conditional-kernel assertion or the conditional product factorization below fails.

For the deterministic trace count, membership of a fixed point \(X_i\) in a closed disk with center \(c\in\mathbb R^2\) and radius \(r\) is equivalent to
\[
 2c^{\mathsf T}X_i+s\ge \lVert X_i\rVert_2^2,
 \qquad s=r^2-\lVert c\rVert_2^2. \tag{1}
\]
Thus every disk trace on \(\{X_1,\ldots,X_n\}\) is a sign pattern induced by \(n\) affine hyperplanes in the three-dimensional parameter \((c_1,c_2,s)\).  Any trace realized by a closed disk can be made strict without changing the trace: if a point is excluded, its distance from the closed disk is positive and the radius may be increased by less than the minimum of these finitely many positive distances; if no point is excluded, any sufficiently small increase works.  Hence equality in (1) creates no additional trace.

Let \(A(n,v)\) be the largest number of open cells cut out by \(n\) affine hyperplanes in \(\mathbb R^v\).  Adding the last hyperplane retains the old cells and splits at most \(A(n-1,v-1)\) of them, because the induced arrangement on that hyperplane has dimension \(v-1\).  Therefore
\[
 A(n,v)\le A(n-1,v)+A(n-1,v-1),
 \qquad A(0,v)=A(n,0)=1. \tag{2}
\]
Induction and Pascal's identity yield
\[
 A(n,3)\le\sum_{j=0}^3\binom nj. \tag{3}
\]
Restricting to genuine disks of a fixed radius, restricting their centers to \(\mathcal B\), and intersecting a trace with the fixed side set \(\{i:X_i\in\mathcal A_t\}\) cannot increase the number of traces.  There are two sides and \(|\mathcal H_n|\) radii by \(\mathrm{def:dyadic\text{-}grid}\), so
\[
 |\mathscr S_n|\le2|\mathcal H_n|\sum_{j=0}^3\binom nj. \tag{4}
\]

If \(\mathscr S_n=\varnothing\), both stochastic conclusions follow from the stipulated zero-maximum convention.  Suppose otherwise.  By \(\mathrm{ass:iid\text{-}sampling}\), regular conditional laws of the observations given \((X_1,\ldots,X_n)\) factor for \(P_X^{\otimes n}\)-almost every score vector; consequently the residuals are conditionally independent.  The partition \(\mathcal A_0=\mathbb R^2\setminus\mathcal A_1\) and sharp assignment give \(T_i=t\) for every index in a side-\(t\) candidate set.  Moreover, the restricted score measure gives full mass to its closed support \(\mathcal X_{t,P}\): its complement is a countable union of rational basic open sets of restricted mass zero.  After removing the two null sets declared in \(\mathrm{ass:gaussian\text{-}errors}\), that assumption therefore gives, for every relevant index and every \(\lambda\in\mathbb R\),
\[
 \mathbb E_P\{e^{\lambda\varepsilon_i}\mid X_1,\ldots,X_n\}
 =\int e^{\lambda u}K_{T_i,P}(X_i,du)
 \le e^{\sigma^2\lambda^2/2}. \tag{5}
\]
The same assumption supplies conditional integrability and zero conditional mean.  No Gaussian equality is used.

For nonempty \(S\in\mathscr S_n\), put
\[
 Z_S=\frac1{\sigma\sqrt{|S|}}\sum_{i\in S}\varepsilon_i.
\]
Conditional independence and (5) imply, for every \(u\in\mathbb R\),
\[
 \mathbb E_P\{e^{uZ_S}\mid X_1,\ldots,X_n\}
 \le\prod_{i\in S}e^{u^2/(2|S|)}=e^{u^2/2}. \tag{6}
\]
Markov's inequality applied to \(e^{zZ_S}\) and then to \(e^{-zZ_S}\) gives
\[
 \Pr_P\{|Z_S|>z\mid X_1,\ldots,X_n\}\le2e^{-z^2/2}.
\]
A union bound over \(S\), which requires no independence among the set-indexed sums, proves
\[
 \Pr_P\left\{\max_{S\in\mathscr S_n}|Z_S|>z
 \ \middle|\ X_1,\ldots,X_n\right\}
 \le2|\mathscr S_n|e^{-z^2/2}. \tag{7}
\]

Finally let \(z_0=\sqrt{2\log\{2|\mathscr S_n|\}}\).  Since \(|\mathscr S_n|\ge1\), \(z_0>1\).  The tail-integral identity for nonnegative random variables, (7), and
\[
 \int_{z_0}^{\infty}e^{-z^2/2}\,dz\le z_0^{-1}e^{-z_0^2/2}
\]
give
\[
 \begin{aligned}
 \mathbb E_P\left[\max_{S\in\mathscr S_n}|Z_S|
 \ \middle|\ X_1,\ldots,X_n\right]
 &\le z_0+2|\mathscr S_n|\int_{z_0}^{\infty}e^{-z^2/2}\,dz\\
 &\le z_0+z_0^{-1}\le z_0+1.
 \end{aligned}
\]
Together with the empty-family case, this is the claimed expectation bound.